\begin{papercard}{nielsen1990}
\cf{Stages} \stagetag{W} (also \stagetag{I}). The original empirical paper on heuristic evaluation: four experiments showing that individual evaluators find only a fraction of the problems, while \emph{aggregates} of independent evaluators do well.

\cf{Definition and heuristics} Heuristic evaluation is informal: looking at an interface and judging good and bad. Instead of guideline collections with about 1{,}000 rules, nine basic heuristics are used: simple and natural dialogue, speak the user's language, minimise memory load, be consistent, provide feedback, clearly marked exits, shortcuts, good error messages, prevent errors \pc{nielsen1990}{1}{We have tried cutting the complexity of the rule base by}.

\cf{Experimental method} Four experiments with non-experts. Evaluators wrote reports on usability problems, and each report was scored by matching to the authors' list of known problems. The list was revised after a first pass because evaluators found problems the authors had missed. Scoring was liberal: credit for mentioning a problem even if incompletely described \pc{nielsen1990}{2}{Each report was then scored for the usability problems that were}. Materials: Teledata screen dumps (37 students, 52 problems); the Mantel one-page specification (77 Computerworld readers, 30 problems); two live voice-response systems (34 students; 48 and 34 problems) \pc{nielsen1990}{2}{Experiment 1 tested the user interlace to the Danish video}.

\cf{Validity of the problem set} Two arguments: the problems are ``obviously'' issues by established knowledge, and at least 34 people worked through each interface. Still, one can never know if all problems were found, so statistics are based on ``known'' problems \pc{nielsen1990}{4}{In spite of these arguments, it is alwa}. A problem counts if it could be expected to bother some users during reasonable use.

\cf{Results} Individual evaluators found on average 51\,\%, 38\,\%, 26\,\% and 20\,\% of problems, so ``heuristic evaluation is difficult'' \pc{nielsen1990}{4}{The most basic resnlt from the four experiments is that}. Performance consistency across two similar systems was weak ($R^2=0.33$) \pc{nielsen1990}{3}{indicates a very weak correlation between the perfor}. Problem finding is not a Guttman scale: poor evaluators sometimes find hard problems and good ones miss easy ones \pc{nielsen1990}{6}{Even poor evaluators can}. False positives were rare and usually found by only one evaluator.

\cf{Aggregation} Monte-Carlo random aggregates show a steep rise from one to about five evaluators and diminishing returns around ten. Five evaluators find about two thirds \pc{nielsen1990}{7}{In general, we would expect}. Evaluators should work \emph{independently} and compare only after writing their reports, to avoid biasing each other \pc{nielsen1990}{7}{For the aggregated evaluation to produce better results than}. Aggregation assumes an authority who can recognise the same problem across reports \pc{nielsen1990}{6}{For tbis aggregation process to work, we have to assume}.

\cf{Conclusions} Use three to five evaluators. HE is cheap and intuitive but gives no direct fixes and is biased by evaluators' mindsets \pc{nielsen1990}{7}{We recommend that heuristic evaluation is done with between}.

\cf{Use for this thesis} Three portable lessons. (1) Single-inspector findings undercount, so a walkthrough by one researcher should be framed as a lower bound, or a second inspector added. (2) \emph{Independent first, merge later} is the basis of a valid aggregation or reliability design. (3) Scoring reports against a ``known problem list'' that is revised after a first pass parallels codebook revision during piloting. Record that the reference list itself changed.
\end{papercard}
